% Generated by tools/textbook/export_listings.js from tools/decks/mastery.js and tools/concepts/registry.json: do not edit.
\begin{chapstatus}{Learning status of this chapter}
\masterychip{now}{UNDERSTAND NOW}\enspace registers, compute, move, jump; registers; functions, call and return; the init, update, draw contract; reading a line of assembly; arguments in a0 to a7; why ra is saved; address, value, load, store; the base instructions; pseudo-instructions; an instruction is a number; tracing by hand.\par\smallskip
\masterychip{use}{USE NOW}\enspace hexadecimal (Chapter~\ref{ch:information}); bit, byte, word (Chapter~\ref{ch:information}); the calling convention (Chapter~\ref{ch:processor}); bit masks (Chapter~\ref{ch:boole1}) and, used without being studied in this course, the assembler directives, the PRG32 tools, the PRG32 runtime calls, the PRG32-QT debugger, the complete program.\par\smallskip
\masterychip{later}{UNDERSTAND LATER}\enspace the stack (Chapter~\ref{ch:processor}); addressing modes (Chapter~\ref{ch:instructions}); instruction formats (Chapter~\ref{ch:instructions}); two's complement (Chapter~\ref{ch:integers}); fixed point (Chapter~\ref{ch:reals}); Boolean algebra (Chapter~\ref{ch:boole1}); the adder (Chapter~\ref{ch:boole2}); from volts to bits (Chapter~\ref{ch:electronics}); latches and flip-flops (Chapter~\ref{ch:sequential}); the processor cycle (Chapter~\ref{ch:instructions}); the Verm1 project (Chapter~\ref{ch:verm1p1}).
\end{chapstatus}
